\documentclass[11pt,a4paper]{article}
\usepackage[T1]{fontenc}
\usepackage[utf8]{inputenc}
\usepackage{lmodern,microtype}
\usepackage[a4paper,margin=25mm,headheight=15pt]{geometry}
\usepackage{amsmath,amssymb,amsthm,mathtools}
\usepackage{booktabs,array,longtable,enumitem}
\usepackage[dvipsnames]{xcolor}
\usepackage{fancyhdr,titlesec,xurl}
\usepackage[unicode,colorlinks=true,linkcolor=MidnightBlue,citecolor=MidnightBlue,urlcolor=MidnightBlue]{hyperref}
\usepackage{bookmark}
\definecolor{ink}{RGB}{27,49,70}
\titleformat{\section}{\large\bfseries\color{ink}}{\thesection}{.65em}{}
\titleformat{\subsection}{\normalsize\bfseries\color{ink}}{\thesubsection}{.65em}{}
\pagestyle{fancy}\fancyhf{}
\fancyhead[L]{\small Poisson layers in feedback transseries}
\fancyhead[R]{\small Research article}
\fancyfoot[C]{\thepage}
\renewcommand{\headrulewidth}{.3pt}
\setlength{\parskip}{3pt}
\setlength{\emergencystretch}{2em}
\setlist{topsep=5pt,itemsep=3pt,leftmargin=1.6em}
\numberwithin{equation}{section}
\newtheorem{theorem}{Theorem}[section]
\newtheorem{proposition}[theorem]{Proposition}
\newtheorem{lemma}[theorem]{Lemma}
\newtheorem{corollary}[theorem]{Corollary}
\theoremstyle{definition}
\newtheorem{definition}[theorem]{Definition}
\newtheorem{example}[theorem]{Example}
\theoremstyle{remark}
\newtheorem{remark}[theorem]{Remark}
\newcommand{\R}{\mathbb R}\newcommand{\N}{\mathbb N}\newcommand{\C}{\mathbb C}
\newcommand{\ee}{\mathrm e}\newcommand{\dd}{\mathrm d}
\newcommand{\E}{\mathbb E}\newcommand{\Prob}{\mathbb P}
\newcommand{\U}{\widehat U}\newcommand{\Acal}{\mathcal A}
\newcommand{\Normal}{\mathcal N}\newcommand{\Pois}{\operatorname{Pois}}
\newcommand{\TV}{d_{\mathrm{TV}}}
\newcommand{\core}{\Phi_M}
\newcommand{\coeff}[2]{[#1^{#2}]}
\hypersetup{pdftitle={Poisson Layers at the Finite-Core Boundary of Exponential-Feedback Transseries},pdfauthor={Research draft prepared with ChatGPT for Vladimir Reshetnikov},pdfsubject={Uniform marked coefficient asymptotics, critical windows, finite-core repair, and joint Gaussian-Poisson limits}}
\begin{document}
\begin{center}
{\LARGE\bfseries\color{ink} Poisson Layers at the Finite-Core Boundary\par of Exponential-Feedback Transseries\par}
\vspace{.65em}
{\large An explicit missing factor, a logarithmically shifted critical window,\par and joint fluctuations of the rare layer and the giant action\par}
\vspace{.8em}
Research article prepared for Vladimir Reshetnikov\\
29 September 2026
\end{center}
\begin{abstract}
Consider the formal equation
\[
 \U(q)=\sum_{j\ge1}q^j\exp(a j^p\U(q)),\qquad a>0,\quad p>1.
\]
A recent ProveIt manuscript proves that a core containing actions
$1,\ldots,M$, together with exactly one larger action, captures the full
$n$th coefficient precisely when $M(p-1)>1$. It leaves the distribution
of the omitted background actions and finite transition corrections open.
We prove a uniform marked-coefficient theorem throughout the larger range
$(M+1)(p-1)>1$. If $w_{n,M}$ is the exact one-tail coefficient and $u_n(s)$
marks each occurrence of action $M+1$ by $s$, then
\[
 u_n(s)=w_{n,M}\exp(s\nu_{n,M})(1+o(1)),\qquad
 \nu_{n,M}=\frac{nr_n}{1+r_n}\ee^{-Mp r_n},
\]
uniformly for real $s$ in compact subsets of $[0,\infty)$. Here
$r_n(1+r_n)^{p-1}\ee^{p r_n}=a n^{p-1}$.
The result is uniform on compact parameter sets with a strict margin, and
survives bounded nonnegative changes to every action inside the core.
Thus one fewer action in the implicit core can be compensated by a single
explicit exponential factor.

At the boundary $p=1+1/M$, the missing factor is exactly
$\exp(r_n^{M+1}/a^M)$ up to relative error tending to zero, not a constant.
A finite Poisson limit instead occurs when
$p=1+1/M+((M+1)\log\log n+\theta)/(M\log n)$; its parameter is
$\ee^{-\theta}/(a^M(M+1)^{M+1})$.
We prove total-variation convergence for bounded intensities, a central
limit theorem and a large-deviation principle for diverging intensities,
relative asymptotics for fixed rare counts, and independence from the
Gaussian fluctuation of the giant action. The proof separates the existing
finite-core method from the new uniform saddle perturbation. Exact
coefficient identities, numerical diagnostics, and further research
questions complete the article.
\end{abstract}
\noindent\textbf{Status and scope.}
This is an AI-assisted research draft with conventional mathematical proofs,
not a Lean-verified development. The results are proposed additions to the
inspected repository work. Global originality and publication priority have
not been established. No general resurgence theorem or signed inverse
conjecture is claimed solved.
\clearpage
\setcounter{tocdepth}{1}\hypersetup{bookmarksdepth=2}
\tableofcontents
\clearpage

\section{The repository question and the contribution}
\label{sec:context}
\subsection{A specific gap in the existing finite-core theory}
The ProveIt repository now keeps its generic transseries work under
\path{Analysis/Transseries/}. Its documentation distinguishes the canonical
transseries volume, focused Lean modules, and recent unmerged research
manuscripts \cite{repo-readme}. The present paper continues
\emph{Finite-Core Universality and Sharp Large Order for Countable
Exponential-Feedback Transseries} \cite{finite-core}, rather than attempting
to replace the entire transseries program by a new survey.

The inspected source is pinned to commit
\begin{center}\small
\path{0c973d8f5e7ef4550f4df1bf486d890037fd9f14}.
\end{center}
Its file is \path{Finite_Core_Universality_Exponential_Feedback/article.tex}
inside the transseries documentation directory, with blob
\path{c2b6faf623003111e9d9e6bbb162a67429f3ce42}.
Its research subsections ``Regularly varying feedback and critical
logarithmic boundaries'' and ``Joint laws for the finite background cloud''
explicitly ask for finite nonzero corrections at a critical cutoff and for
Poisson/Gaussian laws of the small actions. Its necessity argument treats
only one additional occurrence of the first omitted action. A one-occurrence
lower bound proves failure of a cutoff, but does not determine the entire
missing factor.

We sum that layer to relative precision. The result answers the first-omitted-
action part of the cloud question and supplies a moving-power critical
window. It does \emph{not} settle the proposed generalization to arbitrary
regularly varying slopes. It also does not settle the signed quadratic
inverse conjecture in the same source. These boundaries are important:
positivity is used in the proof, and a fixed-parameter theorem cannot be
substituted into a moving-parameter regime without establishing uniformity.

\subsection{What is inherited and what is added}
The positive Lagrange coefficient formula, the giant-action localization,
the strict finite-core cutoff, and the two-stage saddle method originate in
the repository manuscript \cite{finite-core}. We re-establish the forms
needed here, with compact-uniform estimates. They are supporting arguments,
not independent novelty claims.

The additions are the uniform marked ratio, its reduction to a scalar
core-independent intensity, the explicit correction at every reciprocal
cutoff boundary, the shifted Poisson window with its constant, and the joint
rare-layer/giant-action limits. In particular, our coefficient equivalent is
stronger than a logarithmic equivalent: an error $o(1)$ occurs \emph{after}
taking the logarithm of the ratio, not merely after dividing that logarithm
by $n$ or by the diverging intensity.

\subsection{Relation to established mathematics}
Lagrange inversion itself is classical; see Gessel \cite{gessel}. Formal
transseries provide the algebraic context, including the distinction between
formal operations and analytic realization \cite{edgar}. In the present
model, the substitution $q=\ee^{-x}$ produces a one-generator exponential
transseries. Its support is simple; the difficult phenomenon is the growth
and organization of its coefficients.

Poisson layers, condensation, and Gaussian clouds are not new phenomena in
probability. Janson, Jonsson, and Stef\'ansson prove Poisson boundary laws and
joint Gaussian/Poisson limits for superexponentially weighted random trees
\cite{jjs}. Their weights are products of degree weights. Ours contain a
power of the \emph{sum} of action slopes and have a giant of order
$n/\log n$. The logarithmic displacement of the cutoff and the scalar
intensity below are features of this feedback model, not consequences of
identifying the two probability ensembles.

There is also an established notion of mod-Poisson convergence
\cite{modpoisson}. We prove a positive-real-marker exponential approximation.
We do not silently promote it to a relative approximation on an arbitrary
complex marker domain when the intensity diverges. The probabilistic
consequences claimed below are proved from the estimates actually available.

\section{Model, coefficient ensemble, and main theorem}
\label{sec:main}
\subsection{A finite core and a power tail}
Fix an integer $M\ge1$ and put $d=M+1$. To display the universality precisely,
allow a finite nonnegative core
\begin{equation}
 \Phi_M(q,u)=\sum_{j=1}^{M} b_jq^j\ee^{\ell_j u},
 \qquad b_1=1,\quad b_j,\ell_j\ge0.
 \label{eq:core-kernel}
\end{equation}
For a real marker $s\ge0$, let
\begin{equation}
 \U_s(q)=\Phi_M(q,\U_s(q))
       +s q^d\ee^{a d^p\U_s(q)}
       +\sum_{j>d}q^j\ee^{a j^p\U_s(q)},
 \qquad u_n(s)=[q^n]\U_s(q).
 \label{eq:marked-model}
\end{equation}
Every infinite identity involving $\U_s$ is formal. The pure model is
$b_j=1$, $\ell_j=a j^p$, with $s=1$. No evaluation of the divergent infinite
kernel at positive numerical arguments is used.

The parameters may vary with $n$. Uniformity will always mean
\begin{equation}
\begin{gathered}
 a\in[a_-,a_+],\qquad p\in[p_-,p_+],\qquad
 0<a_-\le a_+<\infty,\\
 1<p_-\le p_+<\infty,\qquad d(p_--1)>1,\\
 0\le b_j,\ell_j\le C_0\quad(j\le M),\qquad b_1=1,
 \qquad 0\le s\le S<\infty.
\end{gathered}
\label{eq:uniform-domain}
\end{equation}
Here $M,C_0,S$ and the compact intervals are fixed. This is a strict-margin
condition; no assertion is made at $d(p_--1)=1$.

\subsection{The exact one-tail series}
Let $G=G_M$ and $R=R_M$ be the small analytic branches
\begin{equation}
 G(q)=\Phi_M(q,G(q)),\qquad
 R(q)=\frac{1}{1-\partial_u\Phi_M(q,G(q))}.
 \label{eq:core}
\end{equation}
They satisfy $G(q)=q+O(q^2)$ and $R(q)=1+O(q)$, with nonnegative
Taylor coefficients. Define
\begin{equation}
 W_M(q)=R(q)\sum_{j>M}q^j\ee^{a j^pG(q)},\qquad
 w_{n,M}=[q^n]W_M(q).
 \label{eq:onetail}
\end{equation}
The large-action sum in \eqref{eq:onetail} is formal. The coefficient
$w_{n,M}$ counts exactly one primitive action above $M$, not simply the
coefficient of the finite truncation $G_M$.

\subsection{The scalar intensity}
For $n>0$, let $r=r_n(a,p)>0$ be the unique solution of
\begin{equation}
 r(1+r)^{p-1}\ee^{pr}=a n^{p-1}.
 \label{eq:r}
\end{equation}
Existence and uniqueness follow because the left side increases strictly
from zero to infinity. Set
\begin{equation}
 j_0=\frac{n}{1+r},\qquad k_0=\frac{nr}{1+r},\qquad
 t_0=\ee^{-pr},\qquad
 \boxed{\ \nu_{n,M}=k_0t_0^M=\frac{nr}{1+r}\ee^{-Mp r}.\ }
 \label{eq:nu}
\end{equation}
The intensity depends only on $a,p,n,M$, not on the finite core coefficients.
The root equation is often numerically better evaluated after taking logs.

\begin{theorem}[Uniform Poisson-layer coefficient theorem]
\label{thm:main}
Under \eqref{eq:uniform-domain},
\begin{equation}
 \boxed{\quad
 \sup\left|\frac{u_n(s)}{w_{n,M}\ee^{s\nu_{n,M}}}-1\right|
 \longrightarrow0.
 \quad}
 \label{eq:main}
\end{equation}
The supremum is over all parameters in that domain. In particular,
\begin{equation}
 u_n(1)=w_{n,M}\ee^{\nu_{n,M}}(1+o(1)),\qquad
 u_n(0)=w_{n,M}(1+o(1)).
 \label{eq:repair}
\end{equation}
These are multiplicative equivalents, including when $\nu_{n,M}\to\infty$.
\end{theorem}

\begin{remark}[One fewer implicit core action]
The earlier uncorrected criterion is $M(p-1)>1$. The corrected formula
requires only $(M+1)(p-1)>1$. For $1<p\le2$, the smallest sufficient
uncorrected core has $\lfloor1/(p-1)\rfloor+1$ actions. The present formula
uses $\lfloor1/(p-1)\rfloor$ core actions plus the scalar factor
$\ee^{\nu_{n,M}}$. It does not remove the need to evaluate the one-tail
coefficient or its saddle equivalent.
\end{remark}

\subsection{A directly computable coefficient equivalent}
The one-tail coefficient has a finite-core saddle formula. Let $(z,\tau)$
be the unique saddle in its small-radius asymptotic region:
\begin{equation}
 n=z+a z^p\tau G'(\tau),\qquad
 -\log\tau=apz^{p-1}G(\tau).
 \label{eq:saddle}
\end{equation}
Put $k=n-z$, $\Lambda=a z^p$, and
\begin{align}
 A&=ap(p-1)z^{p-2}G(\tau),\qquad
 B=\Lambda\bigl(\tau G'(\tau)+\tau^2G''(\tau)\bigr),\notag\\
 C&=1+pk/z,\qquad \Delta=C^2-AB,\qquad
 \sigma^2=B/\Delta.
 \label{eq:det}
\end{align}
For large $n$, $\Delta>0$. Define
\begin{equation}
 \Acal_{n,M}=
 \frac{R(\tau)}{\sqrt{\Delta}}
 \exp\bigl(\Lambda G(\tau)-k\log\tau\bigr).
 \label{eq:A}
\end{equation}
We will prove, uniformly on the stated domain,
\begin{equation}
 w_{n,M}\sim\Acal_{n,M},\qquad
 u_n(1)\sim\Acal_{n,M}\ee^{\nu_{n,M}}.
 \label{eq:computable}
\end{equation}
Thus the final formula needs a fixed finite implicit core, one two-variable
saddle, and one increasing scalar equation. It never sums the original
infinite kernel as an analytic function.

\section{Exact formal algebra and probability}
\label{sec:formal}
\begin{lemma}[Positive coefficient identity]
\label{lem:positive}
For nonnegative slopes $\lambda_j$ and amplitudes $c_j$ with $c_1=1$, the
formal equation
$U=\sum_{j\ge1}c_jq^j\ee^{\lambda_jU}$ has a unique solution with
$[q]U=1$. Its coefficients are
\begin{equation}
 [q^n]U=
 \sum_{m=1}^{n}\frac1{m!}
 \sum_{\substack{j_1+\cdots+j_m=n\\j_i\ge1}}
 \left(\prod_{i=1}^{m}c_{j_i}\right)
 \left(\sum_{i=1}^{m}\lambda_{j_i}\right)^{m-1}.
 \label{eq:positive}
\end{equation}
The zeroth power in the one-part term is interpreted as one.
\end{lemma}
\begin{proof}
The coefficient of degree $n$ on the right uses coefficients of $U$ only
through degree $n-1$, which proves triangular existence and uniqueness.
Introduce an auxiliary marker $v$ and solve
$U_v=v\sum c_jq^j\ee^{\lambda_jU_v}$. Formal Lagrange inversion gives
\[
 [v^m]U_v=\frac1m[u^{m-1}]
       \left(\sum_{j\ge1}c_jq^j\ee^{\lambda_j u}\right)^m.
\]
Expanding the exponentials gives \eqref{eq:positive}. Setting $v=1$ is
legitimate coefficientwise, because $m\le n$ for the coefficient of $q^n$.
\end{proof}

For the unmarked model $s=1$, normalize the summands in
\eqref{eq:positive} by $u_n(1)$. This defines a probability measure on
ordered compositions of $n$. Write $K_n$ for the number of parts,
$J_n$ for the largest part, and
\begin{equation}
 E_n=n-K_n+1=1+\sum_i(j_i-1),\qquad
 N_{d,n}=\#\{i:j_i=d\}.
 \label{eq:counts}
\end{equation}
$E_n$ is the total excess plus one; it is not generally the largest part.
Marking in \eqref{eq:marked-model} gives the exact probability generating
function
\begin{equation}
 \E s^{N_{d,n}}=\frac{u_n(s)}{u_n(1)}.
 \label{eq:pgf}
\end{equation}
All distributions in this paper refer to this coefficient ensemble.

Let $\mathcal E_M$ be the event that exactly one part exceeds $M$.
\begin{lemma}[Exact marking interpretation]
\label{lem:onetail}
The branches in \eqref{eq:core} are uniformly analytic near zero on compact
core parameter sets, with nonnegative coefficients, and
\begin{equation}
 \Prob(\mathcal E_M)=w_{n,M}/u_n(1).
 \label{eq:exact-event}
\end{equation}
\end{lemma}
\begin{proof}
On a sufficiently small fixed polydisc, $\partial_u\Phi_M=O(q)$ uniformly.
The analytic implicit function theorem, or a uniform contraction, gives
$G$. Iteration from zero shows positivity. The geometric series defining
$R$ then has nonnegative coefficients as well.

Mark all actions above $M$ by $v$. Differentiating the formal equation at
$v=0$ gives
\[
 (1-\partial_u\Phi_M(q,G))\,\partial_vU_v|_{v=0}
   =\sum_{j>M}q^j\ee^{a j^pG}.
\]
The derivative is $W_M$. Its coefficient counts exactly one marked action
in \eqref{eq:positive}, proving \eqref{eq:exact-event}.
\end{proof}

\section{Compact-uniform finite-core reduction}
\label{sec:uniform}
The following details justify using moving exponents $p=p_n$. They also
isolate the role of positivity. No conclusion about a signed formal inverse
is obtained by this argument.

\subsection{Macroscopic localization}
Let $L=\log n$ and
\begin{equation}
 X_n=(p-1)nL-pn\log L+
 \left(\log a+p\log\frac p{p-1}-(p-1)\right)n.
 \label{eq:X}
\end{equation}
\begin{lemma}[Uniform localization]
\label{lem:macro}
For the marked models with $s\in[0,S]$, $\log u_n(s)=X_n+o(n)$ uniformly.
For any sufficiently small fixed $\varepsilon,\delta>0$, the complement of
\begin{equation}
 \left|E_nL/n-p/(p-1)\right|\le\delta,
 \qquad J_n\ge(1-\varepsilon)E_n
 \label{eq:good}
\end{equation}
has probability at most $\exp(-cn+o(n))$ for a uniform $c>0$. The
probability here is under the $s$-weighted coefficient measure, including
$s=0$. The same logarithmic equivalent holds for any fixed one-tail core.
\end{lemma}
\begin{proof}
Choose uniform constants $D,W$ so that
$\lambda_j\le a j^p+D$ and $c_j\le W$ for all $j$, while $c_1=1$.
A composition with $m$ parts and $r=n-m+1$ has at most $r-1$ nonunit
parts. Convexity and successive transfers to a largest part give
$\sum j_i^p\le r^p+m-1$. Its entire $m$-part row therefore has weight at
most
\begin{equation}
 \binom{n-1}{r-1}W^{r-1}
 \frac{(a r^p+D'm)^{m-1}}{m!}.
 \label{eq:row}
\end{equation}
Conversely, a giant $r>d$ together with unit parts contributes at least
$(a r^p)^{n-r}/(n-r)!$, independent of $s$.

Stirling's formula for the lower configuration at
$r=\lfloor pn/((p-1)L)\rfloor$ gives $X_n+o(n)$, uniformly on compact
parameter sets. We record the upper localization argument to avoid an
unstated uniformity assumption. Rows with $r\ge\eta n$ lose a fixed
multiple of $nL$ for any fixed $\eta>0$. Any remaining row within $O(n)$
of the lower bound must have $r/n\to0$ and
$\log r=L-O(\log L)$; otherwise its slope term is too small.
Thus $D'm=o(r^p)$ and the binomial and amplitude entropies are $o(n)$.
For $\alpha=r/n$, the preceding row estimate is
$(1-\alpha)\{(p-1)L+p\log\alpha+\log a-\log(1-\alpha)+1\}+o(1)$
after division by $n$. Comparing this with the lower bound first gives
$\alpha L=O(\log L)$. Hence $\alpha\log\alpha=o(1)$ and the
quadratic expansion errors also tend to zero. Writing $y=rL/n$, the
row logarithm divided by $n$ becomes
\begin{equation}
 (p-1)L-p\log L+\log a+1+p\log y-(p-1)y+o(1).
 \label{eq:variational}
\end{equation}
The error is uniform in the localized range. Since
$p\log y-(p-1)y$ tends uniformly to minus infinity at both ends of
$(0,\infty)$, comparison with the lower bound now confines $y$ to a
fixed compact interval. This function has its unique maximum at $p/(p-1)$; its maximum gap
outside a fixed neighborhood is bounded below uniformly for $p$ in a
compact interval. Summing at most $n$ rows adds only $O(\log n)$.

If every part is at most $(1-\varepsilon)r$, convexity concentrates the
excess into two extremal parts and bounds the slope sum by
\[
 a r^p\bigl((1-\varepsilon)^p+\varepsilon^p+o(1)\bigr).
\]
The factor in parentheses is uniformly below one. Raising it to
$m-1\sim n$ gives an exponential loss that the $\exp(o(n))$ entropy
cannot offset. This proves the giant estimate. The one-tail lower
configuration and the full coefficient upper bound give the last assertion.
\end{proof}

\subsection{A uniform strict cutoff}
\begin{lemma}[Uniform sufficient core]
\label{lem:cutoff}
For a fixed core size $L_0$ containing all exceptional amplitudes and slopes,
if $L_0(p_--1)>1$, its one-tail coefficient divided by the full coefficient
tends to one uniformly on the compact parameter set.
\end{lemma}
\begin{proof}
On \eqref{eq:good}, the largest part $J$ is unique. If a second part
$i>L_0$ exists, replace $(J,i)$ in their original slots by $(J+i-1,1)$.
The number of parts and their total sum stay fixed; the primitive amplitude
product stays fixed because both large parts are in the unit-amplitude tail.
For any fixed small $\eta>0$, choosing $\varepsilon$ sufficiently small
gives the uniform slope gain
\[
 S'-S\ge a(p-\eta)r^{p-1}(i-1),\qquad r=E_n.
\]
Indeed, the mean-value theorem at the giant gives
$ap(1-\varepsilon)^{p-1}r^{p-1}(i-1)$, and the lost slope of the smaller
part is absorbed by making $\varepsilon$ small. Also
$S'\le a r^p+D'm=a r^p(1+o(1))$. Thus the old weight divided by the
new one is at most
\[
 \exp\left(-\frac{(p-2\eta)(n-r)(i-1)}{r}\right).
\]
Given an image, a chosen slot, and $i$, the preimage is recovered uniquely.
There are at most $n$ slots. Summing over $i>L_0$ bounds the bad weight by
\[
 n\sum_{i>L_0}
 \exp\left(-\frac{(p-2\eta)(n-r)(i-1)}r\right)
 \le n^{1-L_0(p-1)+\epsilon_0+o(1)}
\]
for an arbitrarily small uniform $\epsilon_0>0$, after narrowing the
macroscopic window. The strict compact margin makes this tend to zero.
The exceptional macroscopic set is exponentially negligible. The argument
also works at zero values of exceptional amplitudes, because the transferred
parts are beyond the exceptional core.
\end{proof}

Apply this lemma to the core of size $d=M+1$ in the marked model. Define
\begin{align}
 G_s(q)&=\Phi_M(q,G_s(q))+s q^d\ee^{a d^pG_s(q)},\notag\\
 R_s(q)&=\left(1-\partial_u\Phi_M(q,G_s(q))
               -s a d^p q^d\ee^{a d^pG_s(q)}\right)^{-1},\notag\\
 T_s(q)&=R_s(q)\sum_{j>d}q^j\ee^{a j^pG_s(q)},\qquad
 t_n(s)=[q^n]T_s(q).
 \label{eq:augmented}
\end{align}
Then, uniformly for real $s\in[0,S]$,
\begin{equation}
 u_n(s)=t_n(s)(1+o(1)).
 \label{eq:uniform-reduction}
\end{equation}
At $s=0$, $G_s=G$ and $R_s=R$. The only difference between $T_0$ and
$W_M$ is the fixed $j=d$ summand. That summand is analytic near zero and
has at most exponential coefficient growth. Since $\log w_{n,M}=X_n+o(n)$,
it is negligible, and
\begin{equation}
 t_n(0)\sim w_{n,M}.
 \label{eq:Tzero}
\end{equation}

\section{Two saddles and their uniform perturbation}
\label{sec:saddles}
\subsection{The shrinking Cauchy saddle}
We provide the analytic estimates used later. They involve only a finite
analytic core; the infinite tail is never exchanged with a contour integral.

For $j$ in a compact multiple of $n/\log n$, put $k_j=n-j$ and
$\Lambda_j=a j^p$. Let $t_j$ solve
\begin{equation}
 \Lambda_jt_jG'(t_j)=k_j.
 \label{eq:inner}
\end{equation}
Uniformly, $t_j\sim k_j/\Lambda_j\to0$. Write
$B_j=\Lambda_j(t_jG'(t_j)+t_j^2G''(t_j))\sim n$.
Cauchy's formula gives
\begin{equation}
 [q^{k_j}]R(q)\ee^{\Lambda_jG(q)}
 =\frac{R(t_j)\ee^{\Lambda_jG(t_j)}t_j^{-k_j}}
        {\sqrt{2\pi B_j}}(1+o(1)).
 \label{eq:inner-formula}
\end{equation}
Here and below these assertions hold also for $G_s,R_s$, uniformly in
$s\in[0,S]$.

For completeness, on $q=t_j\ee^{i\theta}$ the linear phase vanishes by
\eqref{eq:inner}. The operator $D=q\,\dd/\dd q$ satisfies
$\Lambda_jD^hG(t_j)=O(n)$ for each fixed positive $h$. On
$|\theta|\le n^{-2/5}$, Taylor expansion therefore gives a Gaussian
quadratic term and an $O(n^{-1/5})$ cubic error. Outside that arc,
nonnegative Taylor coefficients and $[q]G=1$ imply
\[
 \Re\bigl(G(t_j\ee^{i\theta})-G(t_j)\bigr)
 \le t_j(\cos\theta-1).
\]
Since $\Lambda_jt_j\asymp n$, the remaining integral is negligible.
Uniform analyticity controls the amplitude $R$. This proves
\eqref{eq:inner-formula}, including its compact-uniform version.

\subsection{The giant-action saddle}
Define
\[
 f_n(j)=\Lambda_jG(t_j)-(n-j)\log t_j.
\]
Implicit differentiation of \eqref{eq:inner} gives
\begin{align}
 \frac{t_j'}{t_j}&=-\frac{C_j}{B_j},\qquad
 C_j=1+p(n-j)/j,\notag\\
 f_n'(j)&=apj^{p-1}G(t_j)+\log t_j,\notag\\
 f_n''(j)&=ap(p-1)j^{p-2}G(t_j)-C_j^2/B_j.
 \label{eq:derivatives}
\end{align}
Thus, uniformly for $j\asymp n/\log n$,
\begin{equation}
 f_n''(j)=-\frac{pn}{j^2}(1+o(1)),\qquad
 f_n'''(j)=O\left(\frac{(\log n)^3}{n^2}\right).
 \label{eq:curvature}
\end{equation}
At $j=c n/\log n$,
$f_n'(j)/\log n=p/c-(p-1)+o(1)$. The endpoint signs and strict
negative curvature give a unique saddle $z$ in the coarse region, precisely
\eqref{eq:saddle}. Its scales are
\begin{equation}
 z\sim\frac p{p-1}\frac n{\log n},\qquad
 \tau\sim\frac1a\left(\frac{p-1}{p}\right)^p
             \frac{(\log n)^p}{n^{p-1}}.
 \label{eq:scales}
\end{equation}
The equivalences are uniform on compact parameter sets.

Macroscopic localization removes the complementary coarse region with
exponentially small relative error. On $|j-z|\le\sigma\log n$, with
$\sigma^2=B/\Delta\asymp n/(\log n)^2$, Taylor expansion gives
\[
 f_n(j)=f_n(z)-(j-z)^2/(2\sigma^2)+o(1).
\]
The cubic error is $O((\log n)^3/\sqrt n)$. Outside this window, curvature
supplies an $\exp(-c(\log n)^2)$ loss. The prefactor in
\eqref{eq:inner-formula} varies by a relative $o(1)$ on the window.
The lattice Gaussian sum is asymptotic to $\sqrt{2\pi}\sigma$.
Consequently,
\begin{equation}
 w_{n,M}=\Acal_{n,M}(1+o(1)),\qquad
 t_n(s)=\Acal_n(G_s,R_s)(1+o(1)),
 \label{eq:saddle-uniform}
\end{equation}
uniformly. The $\sqrt{2\pi}$ factors from the two saddles cancel.
All assertions about uniqueness here concern the small-radius coarse region,
not unrelated distant solutions of the saddle equations.

\subsection{A normalized Hessian that remains invertible}
The main new comparison can now be made at finite dimension. Define
\begin{equation}
 H_G(j,y)=a j^pG(\ee^y)-(n-j)y.
 \label{eq:H}
\end{equation}
The saddle equations say $\partial_jH_G=\partial_yH_G=0$. Its Hessian at
$(z,\log\tau)$ is
\[
 \begin{pmatrix} A&C\\ C&B\end{pmatrix},\qquad AB-C^2=-\Delta.
\]
It is indefinite, as expected from a Cauchy saddle followed by a maximum
in the action variable. It must not be treated as a positive quadratic form.

Use local coordinates $\xi=(j-z)/z$ and $\eta=y-\log\tau$.
The normalized Hessian is
\begin{equation}
 \frac1n\begin{pmatrix}z^2A&zC\\zC&B\end{pmatrix}
 \longrightarrow
 \begin{pmatrix}p(p-1)&p\\p&1\end{pmatrix}.
 \label{eq:Hnorm}
\end{equation}
The limiting determinant is $-p$, uniformly bounded away from zero.
Derivatives through any fixed order, in these normalized local coordinates,
are $O(n)$ on a fixed small neighborhood. These facts permit ordinary
quantitative implicit-function estimates even though the saddle is indefinite.

\begin{lemma}[First omitted layer perturbs the optimized exponent linearly]
\label{lem:perturb}
Let
\begin{equation}
 \mu=k\tau^M,\qquad k=n-z,
 \label{eq:mu}
\end{equation}
using the $M$-core saddle. Uniformly for $s\in[0,S]$,
\begin{equation}
 \log\frac{\Acal_n(G_s,R_s)}{\Acal_n(G,R)}
 =s\mu+O\left(n\tau^{M+1}+\frac{\mu^2}{n}
                  +\frac{\mu}{n}+\tau^{M+1}\right).
 \label{eq:perturb}
\end{equation}
Every error term tends to zero on \eqref{eq:uniform-domain}.
\end{lemma}
\begin{proof}
Uniform analytic implicit differentiation, or a Taylor expansion of
\eqref{eq:augmented}, yields
\begin{equation}
 G_s(q)-G(q)=s q^d+O(q^{d+1}),\qquad
 R_s(q)-R(q)=O(q^d),\qquad d=M+1.
 \label{eq:jet}
\end{equation}
The same orders hold after any fixed number of Euler derivatives. At the
unperturbed saddle,
\begin{align*}
 H_{G_s}(z,\log\tau)-H_G(z,\log\tau)
 &=s\Lambda\tau^d+O(\Lambda\tau^{d+1})\\
 &=s k\tau^{d-1}+O(n\tau^d),
\end{align*}
because $\Lambda\tau G'(\tau)=k$ and $G'(\tau)=1+O(\tau)$.

The gradient of this perturbation in $(\xi,\eta)$ is $O(\mu)$.
The inverse normalized Hessian is $O(1/n)$ by \eqref{eq:Hnorm}.
The perturbed saddle therefore moves by
\begin{equation}
 (z_s-z)/z=O(\mu/n),\qquad
 \log(\tau_s/\tau)=O(\mu/n).
 \label{eq:shift}
\end{equation}
One can justify these bounds by applying the quantitative inverse-function
argument on a ball of radius $C\mu/n$: the Hessian varies by $o(n)$ there,
and the rescaled Newton map is a contraction. Uniqueness from
\eqref{eq:curvature} identifies the resulting saddle with the one already
used in the coefficient formula. Evaluating at the displaced stationary
point changes the optimized exponent by $O(\mu^2/n)$.

The relative changes of $B$, $\Delta$, and the prefactor are
$O(\mu/n+\tau^d)$, by the derivative bounds and \eqref{eq:jet}.
This proves \eqref{eq:perturb}.

Finally \eqref{eq:scales} gives, up to fixed logarithmic powers,
$n\tau^d=O(n^{1-d(p_--1)}(\log n)^{dp_+})\to0$.
Also $\mu^2/n\le C n\tau^{2M}\to0$, since
$2M(p_--1)>1$: this follows from $d(p_--1)>1$ and $2M\ge M+1$.
The remaining errors are smaller. All margins are strict and uniform.
\end{proof}

\subsection{Why the intensity loses all finite-core dependence}
\begin{lemma}[Additively accurate bare intensity]
\label{lem:bare}
With $\mu$ as in \eqref{eq:mu} and $\nu$ as in \eqref{eq:nu},
\begin{equation}
 \mu-\nu=O(n t_0^{M+1})=o(1)
 \label{eq:bare}
\end{equation}
uniformly under \eqref{eq:uniform-domain}.
\end{lemma}
\begin{proof}
For the bare core $G_0(q)=q$, the saddle equations are solved exactly by
$(j_0,t_0)$ in \eqref{eq:nu}: $a j_0^pt_0=k_0$ and
$-\log t_0=apj_0^{p-1}t_0=pr$.
Every allowed finite core satisfies $G(q)-q=O(q^2)$ uniformly.
The perturbation of its optimized exponent is $O(n t_0)$, and its
normalized gradient is $O(n t_0)$. Applying the same invertible-Hessian
argument as above gives
\[
 (z-j_0)/j_0=O(t_0),\qquad \log(\tau/t_0)=O(t_0).
\]
It follows that
$k\tau^M-k_0t_0^M=O(k_0t_0^{M+1})=O(n t_0^{M+1})$.
The last expression tends to zero by the strict margin.
\end{proof}

\begin{remark}[Two different accuracies]
It is not enough that $\mu/\nu\to1$: both can diverge, and an
$\exp(o(\nu))$ ambiguity would destroy a multiplicative coefficient
formula. Lemma~\ref{lem:bare} gives the stronger \emph{additive} $o(1)$.
By contrast, the bare giant location $j_0$ need not be accurate on the
Gaussian scale. The actual core saddle $z$ remains the center in our joint
limit theorem.
\end{remark}

\begin{proof}[Proof of Theorem~\ref{thm:main}]
Combine the uniform cutoff \eqref{eq:uniform-reduction}, the saddle
formula \eqref{eq:saddle-uniform}, and Lemma~\ref{lem:perturb}:
\[
 u_n(s)=\Acal_{n,M}\exp(s\mu+o(1)).
\]
The errors are uniform even at $s=0$. By Lemma~\ref{lem:bare}, replace
$\mu$ by $\nu$ with additive $o(1)$ in the exponent. Finally
$w_{n,M}\sim\Acal_{n,M}$. Exponentiating a uniform additive $o(1)$ proves
the asserted relative estimate and also \eqref{eq:computable}.
\end{proof}

\section{Every cutoff boundary has an explicit missing factor}
\label{sec:boundary}
\begin{theorem}[Reciprocal-boundary repair]
\label{thm:boundary}
Fix $M\ge1$, $a>0$, and set $p=1+1/M$. Let $r_n>0$ solve
\begin{equation}
 r_n^M(1+r_n)\ee^{(M+1)r_n}=a^M n.
 \label{eq:boundary-r}
\end{equation}
For any fixed allowed finite core,
\begin{equation}
 \boxed{\qquad
 u_n(1)\sim w_{n,M}\exp\left(\frac{r_n^{M+1}}{a^M}\right)
 \sim\Acal_{n,M}\exp\left(\frac{r_n^{M+1}}{a^M}\right).
 \qquad}
 \label{eq:boundary-repair}
\end{equation}
Consequently
\begin{equation}
 \log\frac{u_n(1)}{w_{n,M}}
 \sim\frac{(\log n)^{M+1}}{a^M(M+1)^{M+1}}.
 \label{eq:boundary-log}
\end{equation}
\end{theorem}
\begin{proof}
Here $(M+1)(p-1)=(M+1)/M>1$, so Theorem~\ref{thm:main} applies.
Raising \eqref{eq:r} to the $M$th power gives
\eqref{eq:boundary-r}. That equation gives the exact cancellation
\[
 \nu_{n,M}=\frac{nr_n}{1+r_n}\ee^{-(M+1)r_n}
           =\frac{r_n^{M+1}}{a^M}.
\]
Substitution proves \eqref{eq:boundary-repair}. Since
$r_n\sim\log n/(M+1)$, the logarithmic equivalent follows.
\end{proof}

\begin{example}[The first two boundaries]
For $M=1$, $p=2$, the extra factor beyond the one-action core is
$\exp(r_n^2/a)$, where $r_n(1+r_n)\ee^{2r_n}=an$.
This agrees with, and decomposes, the quadratic correction already found
in \cite{finite-core}; it is a consistency check, not a new quadratic
forward coefficient theorem.

For $M=2$, $p=3/2$, the two-action core misses
$\exp(r_n^3/a^2)$, where $r_n^2(1+r_n)\ee^{3r_n}=a^2n$.
For $M=3$, $p=4/3$, the corresponding factor is
$\exp(r_n^4/a^3)$. The same conclusion holds at every reciprocal boundary.
\end{example}

\subsection{Why the leading logarithm is not a replacement for the root}
Taking logarithms of \eqref{eq:boundary-r}, with $L=\log n$, gives
\[
 (M+1)r_n+M\log r_n+\log(1+r_n)=L+M\log a.
\]
The first two scales are
\begin{equation}
 r_n=\frac{L}{M+1}-\log L+\log(M+1)
             +\frac{M}{M+1}\log a+o(1).
 \label{eq:r-expansion}
\end{equation}
Replacing $r_n$ by $L/(M+1)$ changes $r_n^{M+1}/a^M$ by order
$L^M\log L$, which diverges. Even a relative asymptotic for $r_n$ is
therefore not a relative asymptotic for the missing factor.
The scalar equation is a convenient exact resummation of all logarithmic
terms needed for a coefficient equivalent.

\section{The logarithmically shifted Poisson window}
\label{sec:window}
\begin{theorem}[Critical window and finite survival fraction]
\label{thm:window}
Fix $M\ge1$, $a>0$, and $\theta\in\R$. Let
\begin{equation}
 p_n=1+\frac1M+
 \frac{(M+1)\log\log n+\theta}{M\log n}.
 \label{eq:window}
\end{equation}
Then
\begin{equation}
 \nu_{n,M}\longrightarrow
 \lambda(\theta):=\frac{\ee^{-\theta}}{a^M(M+1)^{M+1}},
 \qquad
 \frac{w_{n,M}}{u_n(1)}\longrightarrow\ee^{-\lambda(\theta)}.
 \label{eq:window-limit}
\end{equation}
The number $N_{M+1,n}$ of first-omitted actions converges in total variation
to $\Pois(\lambda(\theta))$. The conclusions are uniform for $\theta$ in
compact sets and for bounded allowed finite-core perturbations.
\end{theorem}
\begin{proof}
Let $p_0=1+1/M$. On compact parameter sets,
$r_n\sim((p-1)/p)\log n$. From \eqref{eq:r} and \eqref{eq:nu},
\begin{equation}
 \nu_{n,M}
 =a^{-M}n^{1-M(p-1)}r_n^{M+1}
                  (1+r_n)^{M(p-1)-1}.
 \label{eq:nu-power}
\end{equation}
Equivalently,
\begin{align}
 \log\nu_{n,M}
 ={}&[1-M(p-1)]\log n+Mp\log\log n-M\log a\notag\\
 &+Mp\log\frac{p-1}{p}+o(1).
 \label{eq:nu-log}
\end{align}
The $o(1)$ is uniform for the particular moving window, as follows from
the logarithmic root equation and $r_n/\log n\to1/(M+1)$.
Substitute \eqref{eq:window}. The $\log\log n$ terms cancel; terms
involving $(p_n-p_0)\log\log n$ tend to zero. Thus
\[
 \log\nu_{n,M}\to
 -\theta-M\log a-(M+1)\log(M+1).
\]
The coefficient fraction follows from Theorem~\ref{thm:main}; the
probabilistic assertion is proved in Theorem~\ref{thm:probability} below.
\end{proof}

\begin{remark}[The window is not centered at the bare power threshold]
At $p=p_0$ the intensity diverges like a positive power of $\log n$.
Merely setting $p=p_0+\theta/\log n$ still leaves that divergence.
The center must first be moved upward by
$((M+1)/M)\log\log n/\log n$. Only after this displacement does a width
of order $1/\log n$ produce a finite transition profile.
\end{remark}

The constant in \eqref{eq:window-limit} is easily absorbed into a centered
coordinate. With
\[
 \theta=x-M\log a-(M+1)\log(M+1),
\]
the limiting intensity is $\ee^{-x}$ and the surviving one-tail fraction is
$\exp(-\ee^{-x})$. This familiar profile is a consequence of a Poisson
zero-count probability; no extremal-value interpretation is assumed.

\subsection{Finite-size calibration}
Formula \eqref{eq:window} describes a limit, not a high-accuracy numerical
calibration at moderate $n$. The term
$(p_n-p_0)\log\log n$ in \eqref{eq:nu-log} can decay very slowly.
For computation it is preferable to solve
\begin{equation}
 \nu_{n,M}(a,p)=\lambda_*
 \label{eq:calibration}
\end{equation}
using the scalar root \eqref{eq:r}, or to use the exact core intensity
$\mu=k\tau^M$ if the core saddle is already available. This article does
not assert global uniqueness of every finite-$n$ calibration problem.
The numerical routine brackets and checks the chosen local branch.

\section{Poisson, Gaussian, and rare-event consequences}
\label{sec:probability}
Theorem~\ref{thm:main} and \eqref{eq:pgf} give the strong real-marker law
\begin{equation}
 \boxed{\qquad
 \E s^{N_{d,n}}=
 \exp\bigl((s-1)\nu_{n,M}+o(1)\bigr),
 \qquad 0\le s\le S.
 \qquad}
 \label{eq:real-pgf}
\end{equation}
The remainder in the exponent is uniform. In particular it is an additive
$o(1)$ even when the exponent itself diverges.

\begin{theorem}[Distributional and rare-count laws]
\label{thm:probability}
Write $N_n=N_{M+1,n}$ and $\nu_n=\nu_{n,M}$ along any sequence allowed by
\eqref{eq:uniform-domain}.
\begin{enumerate}[label=(\roman*)]
\item If $\nu_n$ stays bounded, then
$\TV(\mathcal L(N_n),\Pois(\nu_n))\to0$.
\item If $\nu_n\to\infty$, then
\[
 \frac{N_n-\nu_n}{\sqrt{\nu_n}}\ \Longrightarrow\ \Normal(0,1).
\]
\item For every regime,
\[
 \Prob(N_n=0)=\exp(-\nu_n)(1+o(1)),\qquad
 \Prob(\mathcal E_M\mid N_n=0)\to1.
\]
\item If $\nu_n\to\infty$, then for each fixed integer $m\ge0$,
\[
 \Prob(N_n=m)\sim\ee^{-\nu_n}\frac{\nu_n^m}{m!}.
\]
\end{enumerate}
\end{theorem}
\begin{proof}
For (i), first pass to any subsequence on which $\nu_n\to\lambda<\infty$.
The pgfs converge on the positive real axis by \eqref{eq:real-pgf}.
For any $S>1$, positivity bounds their absolute values on $|s|\le S$ by
their values at $S$, which remain bounded. Every analytic subsequential
limit therefore equals $\exp(\lambda(s-1))$ by the identity theorem.
Cauchy's coefficient formula gives convergence of every point probability.
The bound at $S>1$ also gives uniform exponentially decreasing tails, so
pointwise convergence upgrades to total variation. Compactness of the
bounded sequence of intensities proves the claim without assuming that
$\nu_n$ itself converges. This includes $\lambda=0$.

For (ii), put $s=\exp(t/\sqrt{\nu_n})$ in \eqref{eq:real-pgf} and subtract
the centering term. For each bounded real $t$,
\[
 \log\E\exp\left(t\frac{N_n-\nu_n}{\sqrt{\nu_n}}\right)
 =\nu_n\left(\ee^{t/\sqrt{\nu_n}}-1-\frac{t}{\sqrt{\nu_n}}\right)+o(1)
 \longrightarrow t^2/2.
\]
The moment-generating-function continuity theorem proves the Gaussian limit.

Setting $s=0$ proves the first part of (iii). For large $n$, configurations
in $\mathcal E_M$ whose unique large part equals $d$ have at most
exponential total coefficient growth and are negligible relative to
$w_{n,M}$. Apart from these, $\mathcal E_M$ implies $N_n=0$.
The identity $u_n(0)\sim w_{n,M}$ therefore proves the conditional claim,
including when the conditioning event is exponentially rare in $\nu_n$.

For (iv), define an analytic polynomial
\[
 F_n(t)=\frac{u_n(t/\nu_n)}{w_{n,M}}.
\]
For every bounded positive real $t$, Theorem~\ref{thm:main} gives
$F_n(t)\to\ee^t$. Positivity gives $|F_n(t)|\le F_n(R)$ on $|t|\le R$.
Normal-family compactness and uniqueness on the positive axis yield
locally uniform analytic convergence to $\ee^t$. Therefore
\[
 [s^m]u_n(s)\sim w_{n,M}\frac{\nu_n^m}{m!}.
\]
Divide by $u_n(1)\sim w_{n,M}\ee^{\nu_n}$.
\end{proof}

\begin{remark}[What is not a consequence here]
For diverging $\nu_n$, we have not proved a total-variation approximation
to $\Pois(\nu_n)$ over \emph{all} counts, a uniform local central limit
theorem, or relative complex-marker control on fixed discs. The Gaussian
limit and the fixed-count relative estimates do not by themselves imply
any of those stronger statements.
\end{remark}

\begin{theorem}[Poisson-rate large deviations]
\label{thm:ldp}
If $\nu_n\to\infty$, the variables $N_n/\nu_n$ satisfy a large-deviation
principle with speed $\nu_n$ and good rate function
\begin{equation}
 I(x)=\begin{cases}x\log x-x+1,&x\ge0,\\+\infty,&x<0,
 \end{cases}\qquad I(0)=1.
 \label{eq:rate}
\end{equation}
\end{theorem}
\begin{proof}
For every real $t$, \eqref{eq:real-pgf} gives
$\nu_n^{-1}\log\E\ee^{tN_n}\to\ee^t-1$.
Chernoff's inequality gives the upper bounds after optimizing over $t$;
for $x>0$, the optimizing value is $t=\log x$ and yields \eqref{eq:rate}.
Positive $t$ also gives exponential tightness.

For the matching lower bound at $x>0$, tilt the measure by $x^{N_n}$.
Its normalized moment generating function follows by taking the ratio of
\eqref{eq:real-pgf} at $x\ee^t$ and at $x$. Chernoff bounds under this
measure show $N_n/\nu_n\to x$ in probability. On a small interval around
$x$, undoing the tilt costs
$\exp(-\nu_n x\log x+o(\nu_n))$, while its normalizer is
$\exp(\nu_n(x-1)+o(\nu_n))$. Let the interval shrink to obtain the lower
bound. At zero the exact relative law for $N_n=0$ supplies the lower bound
$-I(0)=-1$. Standard finite coverings on compact intervals, followed by
exponential tightness, give the full closed-set/open-set formulation.
\end{proof}

\section{The rare layer decouples from the giant action}
\label{sec:joint}
Use the $M$-core saddle $(z,\tau)$ and $\sigma^2=B/\Delta$ from
\eqref{eq:saddle}--\eqref{eq:det}. In particular,
\begin{equation}
 \sigma^2\sim\frac{z^2}{pn}
 \sim\frac{p}{(p-1)^2}\frac{n}{(\log n)^2}.
 \label{eq:sigma}
\end{equation}

\begin{theorem}[Uniform joint transform and independent limits]
\label{thm:joint}
For $s$ in a compact subset of $[0,\infty)$ and real $t$ in a bounded
interval,
\begin{equation}
 \E\left[s^{N_{d,n}}
    \exp\left(t\frac{J_n-z}{\sigma}\right)\right]
 =\exp\left((s-1)\nu_{n,M}+\frac{t^2}{2}+o(1)\right)
 \label{eq:joint}
\end{equation}
uniformly on \eqref{eq:uniform-domain}. Thus:
\begin{enumerate}[label=(\roman*)]
\item If $\nu_n\to\lambda<\infty$, then
\[
 \left(N_{d,n},\frac{J_n-z}{\sigma}\right)
 \Longrightarrow(P,Z),
 \qquad P\sim\Pois(\lambda),\quad Z\sim\Normal(0,1),
\]
with independent limits.
\item If $\nu_n\to\infty$, then
\[
 \left(\frac{N_{d,n}-\nu_n}{\sqrt{\nu_n}},
       \frac{J_n-z}{\sigma}\right)
 \Longrightarrow(Z_1,Z_2),
\]
where $Z_1,Z_2$ are independent standard Gaussians.
\item Conditional on $N_{d,n}=0$, the standardized giant still converges
to $\Normal(0,1)$, even when $\Prob(N_{d,n}=0)\to0$.
\end{enumerate}
\end{theorem}
\begin{proof}
First work in the one-tail representation for the augmented $d$-core. The
largest action is its unique tail index $j$. Multiplying its outer saddle
sum by $\exp(t(j-z)/\sigma)$ tilts a lattice Gaussian. Taylor expansion
on the same window as before gives the factor
\[
 \exp\left(t\frac{z_s-z}{\sigma}
              +\frac{t^2\sigma_s^2}{2\sigma^2}\right)(1+o(1)).
\]
By \eqref{eq:shift} and \eqref{eq:sigma},
\[
 \frac{z_s-z}{\sigma}=O(\mu/\sqrt n)=o(1),\qquad
 \sigma_s/\sigma\to1.
\]
Thus the factor is $\exp(t^2/2)(1+o(1))$, uniformly in both markers.

We justify transferring this tilted estimate to the full positive
coefficient model, rather than appealing to convergence in probability of
the cutoff event alone. The giant tilt changes any weight by at most
$\exp(O(n/\sigma))=\exp(O(\sqrt n\log n))$, which cannot overcome the
macroscopic exponential loss in Lemma~\ref{lem:macro}. In the transfer map
of Lemma~\ref{lem:cutoff}, the giant changes by $i-1$. The extra ratio is
at most $\exp(|t|(i-1)/\sigma)$. Its exponent per transferred unit is
$O(\log n/\sqrt n)$, negligible against the existing suppression of order
$\log n$. The strict cutoff estimate therefore remains uniform for the
positively tilted weights. This proves the required relative reduction.

Combining the tilted saddle factor with Theorem~\ref{thm:main} proves
\eqref{eq:joint}. Substituting $s=\ee^v$ in the bounded-intensity case
identifies the product limiting mgf. In the diverging case substitute
$s=\exp(v/\sqrt{\nu_n})$ and center as in
Theorem~\ref{thm:probability}; the limit is
$\exp((v^2+t^2)/2)$. Finally set $s=0$ in \eqref{eq:joint} and divide by
the zero-count probability to obtain the conditional Gaussian mgf.
\end{proof}

\begin{remark}[The center must retain the core]
The leading estimate $z\sim pn/((p-1)\log n)$ does not determine a valid
Gaussian center. Even ordinary next logarithmic corrections can exceed
$\sigma$ by an unbounded factor. Likewise the core-to-bare shift can be
large relative to $\sigma$. Core independence of the scalar intensity
$\nu$ does not imply core independence of the finely centered giant.
\end{remark}

\section{Computation and checks}
\label{sec:checks}
\subsection{Two independent exact coefficient algorithms}
For the power slope $\lambda_j=a j^p$, set
$E_{j,m}=[q^m]\exp(\lambda_jU(q))$. The triangular recurrence is
\begin{equation}
 u_n=\sum_{j=1}^{n}c_jE_{j,n-j},\qquad
 E_{j,0}=1,\qquad
 E_{j,m}=\frac{\lambda_j}{m}
       \sum_{k=1}^{m}k u_kE_{j,m-k}.
 \label{eq:recurrence}
\end{equation}
For a finite core, compute $G$ by the same recurrence with $c_j=0$ above
$M$. Then compute the coefficients of
$P=\partial_u\Phi_M(q,G)$, $R=(1-P)^{-1}$, and the convolution
\eqref{eq:onetail}. All these calculations are finite to each requested
order.

The independent check enumerates integer partitions. For a multiplicity
vector $(m_j)$ with $\sum j m_j=n$, its contribution is
\begin{equation}
 \frac{\prod_j c_j^{m_j}}{\prod_jm_j!}
       \left(\sum_j\lambda_jm_j\right)^{\sum_jm_j-1}.
 \label{eq:partition}
\end{equation}
To check the one-tail series, retain only vectors with
$\sum_{j>M}m_j=1$. This is not a recurrence enforcing the same identity;
it is a separate finite enumeration based on the formal coefficient formula.

The supplied program uses exact rational arithmetic for $a=1$, $p=2,3$,
and degrees through $16$. It compares the triangular recurrence with the
partition sum for markers on action $2$ or $3$, at marker values $0,1,2$;
it separately compares the one-tail recurrence with the restricted partition
sum for $M=1,2,3$. All \textbf{288 exact assertions passed}. A further
32 comparisons of the log-domain numerical recurrence with the exact
coefficients had maximum absolute logarithmic discrepancy
$7.11\times10^{-15}$ in the recorded run.

For example, at $a=1,p=2$, the first eight forward coefficients are
\[
 1,\ 2,\ \frac{15}{2},\ \frac{107}{3},\ \frac{1555}{8},\
 \frac{17362}{15},\ \frac{5288311}{720},\ \frac{15414767}{315}.
\]
These checks detect normalization and marking mistakes. They do not prove
an asymptotic theorem or certify a floating-point saddle.

\subsection{Coefficient diagnostics with finite-size calibration}
For each row of Table~\ref{tab:diagnostics}, the exponent $p$ was chosen
numerically so that the core intensity $\mu=k\tau^M$ equals $1/2$.
The table compares exact-event coefficients computed by the positive
log-domain recurrence, not by the asymptotic formula itself. Values are
ordinary floating-point diagnostics, not outward-rounded enclosures.

\begin{table}[htbp]
\centering\small
\caption{Calibrated coefficient diagnostics, $a=1$. The prediction is
$\ee^{-\nu}$ with the bare scalar intensity.}
\label{tab:diagnostics}
\begin{tabular}{rrrrrr}
\toprule
$M$&$n$&$p$&$w_{n,M}/u_n(1)$&$\Prob(N_{d,n}=0)$&$\ee^{-\nu}$\\
\midrule
1&80 &2.562513&0.592485&0.595906&0.607610\\
1&160&2.554385&0.600168&0.601666&0.607037\\
1&320&2.539670&0.603573&0.604266&0.606770\\
2&80 &1.707730&0.377341&0.449096&0.596771\\
2&160&1.728827&0.471136&0.508385&0.596605\\
2&320&1.739877&0.520370&0.541980&0.598192\\
\bottomrule
\end{tabular}
\end{table}

The slower convergence for $M=2$ is visible and should not be concealed.
The normalized coefficient ratio $u_n/(\Acal_{n,M}\ee^{\nu})$ is
$1.02495,1.01113,1.00514$ in the first three rows and
$1.58321,1.26638,1.14950$ in the last three. The theory asserts a limit,
not that a minimal corrected core is already highly accurate at $n=80$.
The auxiliary CSV files retain the unrounded values and residuals.

\subsection{A warning from the explicit window}
The second diagnostic table evaluates the scalar/core saddles at very large
$n$, using \eqref{eq:window} with
$\theta=-(M+1)\log(M+1)$ and $a=1$. The limiting intensity is then one.
Even at these large values, the finite-size intensity need not be close to
one: for $M=1$, the computed core intensity is approximately $1.24932$ at
$n=10^6$, $1.33102$ at $10^{12}$, and $1.23906$ at $10^{48}$.
This is consistent with slowly vanishing logarithmic corrections; it is not
evidence of a monotone convergence theorem. The full diagnostic file records
$M=1,2,3$. For numerical work the root-defined $\nu$ or the calibrated
$\mu$ should be used instead of the limiting constant alone.

\subsection{Reproducibility and complexity}
The bundle contains \path{verification/verify.py}, recorded JSON/CSV outputs,
and a run log. The exact checks require Python's standard rational arithmetic;
the same script also imports NumPy and SciPy for the numerical part. The
recorded environment used Python 3.13.5, NumPy 2.3.5, and SciPy 1.17.0.
The direct triangle takes $O(N^3)$ arithmetic operations and $O(N^2)$
storage; this is not an exact-arithmetic bit-complexity estimate.

The saddle routine checks the implicit residual, the regular core branch,
and the determinant sign. These checks are necessary but not interval
certificates. A small floating residual alone is not a proof that a numerical
root has been enclosed. No external service is used by the verification
program, and no repository file was modified in producing this package.

\section{Further research questions}
\label{sec:future}
\subsection{A full local Poisson approximation at diverging intensity}
The real-marker theorem gives a central limit theorem, fixed rare-count
equivalents, and large deviations. Prove a uniform local approximation for
$m=\nu+O(\sqrt\nu\log\nu)$ and then total-variation convergence to
$\Pois(\nu)$ when $\nu\to\infty$. A promising approach is to extract the
marker coefficient within the two-saddle calculation, controlling the
conditioning correction $\nu^2/n$ directly. This would be stronger than
merely invoking the Gaussian weak limit.

\subsection{The next omitted layer and interaction corrections}
Our compact margin makes $n\tau^{M+1}$ and $\mu^2/n$ vanish. At the next
boundary these terms can survive. Determine the joint law of actions
$M+1$ and $M+2$ and the first interaction correction to the exponential
factor. The normalized Hessian in \eqref{eq:Hnorm} gives a concrete starting
point: the optimized second-order correction is a quadratic form in the
marked perturbation gradients. It need not have the sign of a covariance
matrix, because the underlying two-variable saddle is indefinite.

\subsection{Fixed logarithmic corrections to the power tail}
For slopes $a j^p(\log j)^\gamma$, the repository predicts a new cutoff
behavior at $M(p-1)=1$. Extend the compact-uniform transfer and saddle
estimates to a differentiably regularly varying tail. The delicate target
is the finite correction at $\gamma=p$, including its constant and
Poisson parameter. The present moving-$p$ theorem does not prove this
fixed-tail statement, although its perturbation mechanism suggests what to
measure.

\subsection{All finite cloud coordinates and the correct centers}
Study the joint vector of counts of actions $2,\ldots,M+1$, together with
the giant. Some counts diverge on different scales, and leading equivalents
for their means may be too inaccurate for central limits. Introduce several
markers in the finite core, solve the corresponding marked saddle equations,
and compute the complete covariance and centering corrections. The present
paper resolves the first omitted coordinate, not this whole multiscale vector.

\subsection{The signed quadratic inverse conjecture}
The source \cite{finite-core} conjectures a sharp equivalent for the formal
inverse at $p=2$. Its exact inverse one-tail response is known, but positivity
is lost. A signed finite-core reduction must control cancellations before a
Poisson-type exponent can be trusted. Our transfer estimates cannot simply
be applied with absolute values: doing so may swamp the proposed inverse
coefficient. An analytic-core regrouping of the inverse Lagrange formula is
still a focused and worthwhile next problem.

\subsection{Quantified error bounds and certified coefficient budgets}
Replace the compact-uniform $o(1)$ terms by explicit bounds with numerical
onset values. The perturbation estimates already identify two obstruction
scales, $n\tau^{M+1}$ and $\nu^2/n$; the cutoff and contour bounds must
supply computable constants as well. Combine these with interval evaluation
of the finite core and a certified saddle enclosure. The objective is an
actual coefficient interval, not merely a residual-based diagnostic.

\subsection{A growing core near the linear-feedback boundary}
As $p\downarrow1$, the required core size grows. A theorem allowing
$M=M_n\to\infty$ would require uniform analytic radii and uniform control of
all finite-core derivatives. The present proof fixes $M$ throughout. The
interaction between a growing core and the scale $(p-1)\log n$ is not
captured by taking a formal limit in \eqref{eq:main}.

\subsection{Formalization and analytic realization}
A useful Lean development can begin with the finite marking identities,
the scalar monotonicity in \eqref{eq:r}, the exact boundary cancellation,
and a reusable parameter-uniform saddle-perturbation lemma. The later
probability and contour arguments should be separate modules with explicit
hypotheses. In parallel, investigate whether these coefficient-layer laws
control optimally truncated analytic realizations on the negative ray.
A least formal term or a coefficient probability law is not automatically
an analytic remainder theorem.

\section{Conclusion}
The strict finite-core cutoff leaves a quantitatively structured remainder.
In the range where one more core action would suffice, the entire first
omitted layer contributes a single exponential factor with an explicit
scalar intensity. That intensity is insensitive, to additive $o(1)$, to
bounded nonnegative changes inside the core.

At every reciprocal boundary the correction grows like the exponential of a
power of $\log n$. A finite nonzero survival fraction emerges only after a
$\log\log n/\log n$ displacement of the exponent. The same coefficient
theorem yields Poisson limits, Gaussian fluctuations, exponentially small
fixed-count probabilities, and a joint limit in which the rare layer and
the giant fluctuation decouple.

These results extend specific repository questions without confusing formal
coefficient analysis with analytic summability, signed reversion, or a
machine-checked development. The unresolved extensions above are stated as
separate mathematical tasks rather than conclusions inferred by analogy.

\appendix
\section{Proof-dependency and scope audit}
\label{app:audit}
\begin{longtable}{>{\raggedright\arraybackslash}p{.25\textwidth}
                  >{\raggedright\arraybackslash}p{.42\textwidth}
                  >{\raggedright\arraybackslash}p{.24\textwidth}}
\toprule
Statement & Main inputs & Status here\\
\midrule\endhead
Positive coefficients and one-tail identity & Formal Lagrange inversion;
finite marking & Proved; classical tools\\[5pt]
Compact-uniform sufficient core & Positive transfer, preimage count,
uniform convex localization & Re-established from repository method\\[5pt]
Uniform finite-core saddle & Cauchy Gaussian estimate, action saddle,
lattice summation & Re-established in needed scope\\[5pt]
Marked exponential repair & Uniform core jet and invertible normalized
indefinite Hessian & Main extension\\[5pt]
Core-independent intensity & Additive bare-saddle comparison & Main extension\\[5pt]
Boundary correction and shifted window & Exact scalar cancellations and
uniform theorem & Main extension\\[5pt]
Bounded-intensity total variation & Positive pgfs, analytic compactness,
exponential tightness & Proved\\[5pt]
Diverging-intensity CLT and fixed counts & Real-marker theorem; scaled
analytic compactness & Proved\\[5pt]
Giant/layer independence & Positive tilt, uniform cutoff, shifted Gaussian
saddle & Proved\\[5pt]
Signed inverse; general regular variation; growing core & Additional
cancellation or uniformity arguments required & Not claimed\\
\bottomrule
\end{longtable}
The full power tail is exact outside the specified finite core. A merely
asymptotically equivalent tail, unbounded signed amplitudes, a growing $M$,
or approach to the excluded lower boundary require additional work.
All asymptotic assertions are as $n\to\infty$. Computations in this package
are checks of formulas and diagnostics, not additional hypotheses in a proof.

\clearpage
\section{Notation and an implementation checklist}
\begin{longtable}{>{\raggedright\arraybackslash}p{.2\textwidth}
                  >{\raggedright\arraybackslash}p{.71\textwidth}}
\toprule
Symbol & Meaning\\\midrule\endhead
$M$, $d=M+1$ & Core size and first omitted primitive action\\[3pt]
$\U_s$, $u_n(s)$ & Formal marked feedback solution and its $n$th coefficient\\[3pt]
$G,R$ & Analytic finite-core solution and linear response\\[3pt]
$W_M,w_{n,M}$ & Formal exactly-one-tail series and coefficient\\[3pt]
$G_s,R_s,T_s$ & Augmented $d$-core and its exactly-one-tail response\\[3pt]
$z,\tau,k$ & Actual $M$-core saddle action, radius, and $n-z$\\[3pt]
$r,j_0,t_0,k_0$ & Scalar bare root and bare saddle quantities\\[3pt]
$\mu=k\tau^M$ & Core-dependent intensity used in the perturbation proof\\[3pt]
$\nu=k_0t_0^M$ & Explicit core-independent intensity; $\mu-\nu=o(1)$\\[3pt]
$A,B,C,\Delta$ & Hessian entries and positive saddle determinant factor\\[3pt]
$\Acal_{n,M}$ & Finite-core saddle equivalent for $w_{n,M}$\\[3pt]
$N_{d,n}$ & Number of primitive action-$d$ occurrences in the coefficient ensemble\\[3pt]
$J_n,\sigma$ & Largest action and its Gaussian scale\\
\bottomrule
\end{longtable}

For numerical use, first choose a core with $(M+1)(p-1)>1$ and a positive
margin. Solve \eqref{eq:r} in logarithms, solve the finite-core saddle on
its regular small branch, and evaluate \eqref{eq:A} in logarithms. Add
$\nu$ to its logarithm, not a leading asymptotic replacement for $\nu$.
Check the branch and residual, but do not interpret those checks as an
interval enclosure. Increase the core or use a more accurate correction
when the finite-size residual scales are not small.

\begingroup\small
\begin{thebibliography}{9}
\setlength{\itemsep}{3pt}\setlength{\parskip}{0pt}
\bibitem{repo-readme}
V. Reshetnikov, \emph{ProveIt: transseries documentation and source inventory},
repository documentation inspected 29 September 2026.
\url{https://github.com/VladimirReshetnikov/ProveIt/tree/0c973d8f5e7ef4550f4df1bf486d890037fd9f14/Analysis/Transseries}.

\bibitem{finite-core}
\emph{Finite-Core Universality and Sharp Large Order for Countable
Exponential-Feedback Transseries}, AI-assisted research draft prepared for
V. Reshetnikov, 29 September 2026. ProveIt source at the pinned commit.
\url{https://github.com/VladimirReshetnikov/ProveIt/blob/0c973d8f5e7ef4550f4df1bf486d890037fd9f14/Analysis/Transseries/docs/series-and-transseries/Finite_Core_Universality_Exponential_Feedback/article.tex}.
In particular, the strict cutoff and saddle theorems, the one-omitted-action
necessity proof, and the research subsections on logarithmic boundaries and
finite background clouds.

\bibitem{gessel}
I. M. Gessel, \emph{Lagrange inversion}, Journal of Combinatorial Theory,
Series A \textbf{144} (2016), 212--249.
\url{https://arxiv.org/abs/1609.05988}.

\bibitem{edgar}
G. A. Edgar, \emph{Transseries for beginners}, Real Analysis Exchange
\textbf{35} (2009/2010), 253--310.
\url{https://arxiv.org/abs/0801.4877}.

\bibitem{jjs}
S. Janson, T. Jonsson, and S. \"O. Stef\'ansson,
\emph{Random trees with superexponential branching weights},
Journal of Physics A: Mathematical and Theoretical \textbf{44} (2011), 485002.
\url{https://arxiv.org/abs/1104.2810}.
Theorems 2.5--2.6 and Remark 2.7 provide the relevant comparison with
Poisson and Gaussian degree-count limits in a different ensemble.

\bibitem{modpoisson}
E. Kowalski and A. Nikeghbali, \emph{Mod-Poisson convergence in probability
and number theory}, International Mathematics Research Notices
\textbf{2010}, no. 18, 3549--3587.
\url{https://arxiv.org/abs/0905.0318}.
\end{thebibliography}
\endgroup
\end{document}
